\documentclass[12pt, letterpaper]{article}

% --- Paquetes básicos ---
\usepackage[spanish, es-tabla]{babel}
\usepackage[utf8]{inputenc}
\usepackage[T1]{fontenc}
\usepackage{newpxtext,newpxmath} % Fuente estilo Palatino
\usepackage[margin=3cm]{geometry}
\usepackage{titlesec}
\usepackage[autostyle]{csquotes}

% --- Para dibujar la V de Gowin ---
\usepackage{tikz}
\usetikzlibrary{positioning,babel,calc}

% --- Configuración de la sección (1. IDEA DE...) ---
\titleformat{\section}
  {\normalfont\large\scshape\centering} 
  {\thesection.}
  {0.5em}
  {}

\begin{document}

% --- Encabezado Manual ---
\begin{center}
    \vspace{0.8cm}
    {\Large \bfseries A STUDY OF PROBABILISTIC PASSWORD MODELS} \\
    \vspace{0.4cm}
\end{center}

\vspace{1.5cm}

\noindent
\textbf{Nivel descriptivo} \newline

\noindent
Título: A Study of Probabilistic Password Models \newline

\noindent
Autor: Jerry Ma, Weining Yang, Min Luo y Ninghui Li \newline

\noindent
Año: 2014 \newline

\noindent
Nombre del tema: Estudio de modelos probabilísticos para la predicción de contraseñas. \newline

\noindent
Clasificación: Metodológica; los autores comparan distintos modelos probabilísticos para estimar la fuerza y la probabilidad de aparición de contraseñas dentro de grandes corpus.\newline

\noindent
Resumen en una oración: El artículo muestra que los modelos probabilísticos bien diseñados pueden predecir contraseñas con bastante más precisión que enfoques más simples. \newline

\noindent
Argumento central: Los autores sostienen que, si se usan adecuadamente técnicas de normalización, suavizado y modelado estadístico, es posible construir modelos de ataque más eficaces para contraseñas.\newline

\noindent
Problemas con el argumento: Una dificultad importante es que la eficacia del modelo depende fuertemente del conjunto de datos usado y de las decisiones técnicas tomadas durante el entrenamiento, por lo que los resultados no siempre se trasladan bien a otros contextos. Además, los modelos más complejos pueden ser costosos de calcular y difíciles de comparar con otros enfoques más simples.
\newpage 

\noindent
\textbf{Nivel analítico} \newline

\noindent
Conexión con mi proyecto: Este trabajo se relaciona directamente con mi proyecto porque propone una manera de estudiar las contraseñas como objetos probabilísticos, lo cual sirve para comparar la frecuencia esperada de aparición con otras medidas de seguridad o entropía.\newline

\noindent
Lo que NO dice el autor: Los autores no analizan en profundidad cómo cambia la estructura de las contraseñas según el contexto social o la plataforma donde fueron creadas, ni estudian cómo influye el uso de reglas humanas más creativas en el proceso de generación. \newline

\noindent
Contraste con otra fuente: A diferencia de estudios como el de Bonneau, que se enfocan más en la dificultad real de adivinar contraseñas y en criticar métricas tradicionales, este artículo pone el énfasis en construir modelos de ataque más precisos desde el punto de vista computacional.   \newline 

\noindent
Evaluación de aplicabilidad: 4. Aunque no resuelve por sí solo el problema de medir la entropía de contraseñas en todos los casos, sí ofrece herramientas estadísticas útiles para construir comparaciones más formales entre conjuntos de datos. Además, su enfoque es bastante reproducible y puede adaptarse a análisis posteriores. \newline

\noindent
Preguntas que le haría al autor: ¿Cómo cambiaría el desempeño del modelo si se normalizaran previamente patrones de teclado, caracteres Unicode o variaciones lingüísticas antes de entrenarlo? \newline

\noindent
Resumen argumentativo: Los autores intentan resolver el problema de predecir contraseñas mediante modelos probabilísticos más refinados. Su propuesta es valiosa porque muestra que la precisión del análisis depende mucho de la forma en que se preparan y suavizan los datos. Para el proyecto, esto sirve como base metodológica para evitar sesgos al calcular probabilidades o frecuencias, especialmente cuando se desea comparar contraseñas humanas con patrones generados automáticamente.

\end{document}
